\section{X² + Y² = p}

\begin{proposicao}
Seja $\mathcal{M}$ um conjunto não-vazio de inteiros Gaussianos. Se $\mathcal{M}$ é fechado com respeito à subtração
e satisfaz a condição ``se $a$ pertence a $\mathcal{M}$, então todos os múltiplos de $a$ pertencem a $\mathcal{M}$'',   
então $\mathcal{M}$ consiste em todos os múltiplos de algum inteiro Gaussiano $d$, unicamente determinado a menos de um 
fator invertível.
\end{proposicao}

\begin{proof}
Se $\mathcal{M}=\{0\}$, a proposição é verdadeira com $d=0$.  
Senão, tomemos em $\mathcal{M}$ um elemento $d$ de menor norma $>0$.  
Se $a \in \mathcal{M}$, podemos aplicar a divisão euclidiana em $\mathbb{Z}[i]$ (Lema 7.1)
e escrever $a = dq + r$ com $N(r) \leq \tfrac{1}{2}N(d)$.  
Então, $r = a - dq \in \mathcal{M}$ e portanto $r=0$,ou seja, $a$ é um múltiplo de $d$.
Quanto à unicidade, se $d'$ possui a mesma propriedade, então $d$ e $d'$ dividem um ao outro, 
e portanto são associados.  
\end{proof}

Assim como nas Aulas 4 e 6, podemos aplicar a Proposição 8.1 ao conjunto de todas as combinações lineares
$ax + by + \ldots + cz$ onde $a,b, \ldots, c$ são inteiros Gaussianos dados e $x,y, \ldots, z$ são inteiros 
Gaussianos arbitrários, e assim definir o $m.d.c. (a,b, \ldots, c)$. Este será unicamente determinado se 
prescrevermos que ele deve ser \emph{normalizado} (no sentido da Aula 7). Se ele for 1, diremos que 
$a,b, \ldots, c$ são mutuamente relativamente primos. Podemos, então, repetir todos os conteúdos da Aula 4 
(\emph{e.g} a versão análoga da Proposição 4.8 seria: se um primo Gaussiano divide um produto de certos inteiros
Gaussianos, então ele divide pelo menos um dos fatores), exceto que na demonstração do Teorema 4.9 foi utilizado 
indução sobre um inteiro $a$, enquanto que agora usaremos indução sobre a norma $N(a)$. A conclusão é :

\begin{proposicao}
Todo inteiro Gaussiano não-nulo pode ser escrito de modo
``essencialmente único'' como um produto de um invertível e de primos
Gaussianos.
\end{proposicao}

Aqui, as palavras ``essencialmente único'' têm o seguinte sentido.
Sejam
\[
a = u q_1 \cdots q_r = u' q'_1 \cdots q'_s
\]
dois produtos do tipo desejado para algum $a \neq 0$, em que $u$ e $u'$
são invertíveis e $q_j$ e $q'_k$ são primos Gaussianos. Então, o teorema
deve ser entendido por dizer que $r = s$ e que os $q'_k$ podem ser
ordenados de modo que $q'_j$ seja um associado de $q_j$ para
$1 \le j \le r$; se $a$ for invertível, então $r = 0$.
Se escrevermos que os fatores primos de $a$ sejam ``normalizados'',
então o produto é unicamente determinado a menos da ordem dos fatores.

Inteiros ordinários também são inteiros Gaussianos; para obtermos a
decomposição deles em primos Gaussianos, basta fazermos isso para os
primos ``ordinários''.

\begin{proposicao}
Se $p$ é um primo racional ímpar, então ele é um primo Gaussiano ou a
norma de um primo Gaussiano $q$; neste último caso,
$p = q\bar{q}$, e $q$ e $\bar{q}$ não são associados, e os únicos
divisores de $p$ são $q$, $\bar{q}$ e os associados deles.
\end{proposicao}

\begin{proof}
Escrevamos $p = u q_1 \cdots q_r$ como na Proposição 8.2.
Tomando a norma, encontramos $p^2 = N(q_1) N(q_2) \cdots N(q_r)$.
Se um dos $N(q_j)$ for $p^2$, então $r = 1$, $p = u q_j$, e $p$ seria um
primo Gaussiano. Caso contrário, cada $N(q_j)$ é $p$, e podemos escrever
$p = N(q) = q \bar{q}$, com $q$ e $\bar{q}$ primos Gaussianos. Digamos 
que $q = x + iy$; se $\bar{q}$ fosse um associado de $q$, ele seria 
$\pm q$ ou $\pm iq$; disso, teríamos $y = 0$ e $p = x^2$, ou $x = 0$ e $p = y^2$, 
ou $y = \pm x$ e $p = 2x^2$, o que não pode acontecer, pois $p$ é um primo ímpar.
\end{proof}

Para $p = 2$, temos a decomposição
\[
2 = N(1 + i) = (1 + i)(1 - i) = i^3(1 + i)^2,
\]
ou seja, 2 possui um único fator primo normalizado $1 + i$.

\begin{proposicao}
Se $p$ é um primo racional ímpar, então $p$ é um primo Gaussiano ou a
norma de um primo Gaussiano de acordo quando ele deixar resto $3$ ou $1$
na divisão por $4$.
\end{proposicao}

\begin{proof}
Se $p$ for a norma de $x + iy$, então $p = x^2 + y^2$,
em que um dos inteiros $x$ ou $y$ deve ser ímpar e o outro par.
Logo, um dos quadrados $x^2$ ou $y^2$ deve deixar resto $1$ na divisão
por $4$, enquanto o outro deixa resto $0$.

Reciprocamente, se $p - 1$ é um múltiplo de $4$, existe um inteiro
ordinário $x$ tal que $x^2 + 1$ é um múltiplo de $p$ (não é difícil ver
isso de várias maneiras; e.g., Corolário 17.6). Como
$x^2 + 1 = (x + i)(x - i),$
isto, se $p$ fosse um primo Gaussiano, resultaria em $p$ dividir
$x + i$ ou $x - i$. Simplesmente, isso não pode acontecer. 
\end{proof}

\begin{corolario}
Cada primo Gaussiano é $\pm 1 \pm i$, ou um associado de um primo
racional que deixa resto $3$ na divisão por $4$, ou ainda sua norma é
um primo racional que deixa resto $1$ na divisão por $4$.    
\end{corolario}

\begin{proof}
Cada primo Gaussiano $q$ deve dividir algum fator primo racional $p$
de sua norma $q\bar{q}$; aplicando a Proposição~8.4 se $p$ for ímpar,
e as observações anteriores se $p = 2$, obtemos o desejado.
\end{proof}

\begin{corolario}[Fermat]
Um primo racional pode ser escrito como uma soma de dois quadrados se,
e somente se, ele é igual a $2$ ou deixa resto $1$ na divisão por $4$.
\end{corolario}

\begin{proof}
Se $p = x^2 + y^2$, então $p$ não pode ser um primo Gaussiano pois ele
possui divisores $x \pm iy$. A recíproca é dada pela
Proposição~8.4.
\end{proof}

Vale ressaltar que este famoso resultado sobre inteiros racionais foi
mostrado utilizando-se de um anel maior, o anel dos inteiros
Gaussianos.
